\begin{proof}[of
\Cref{prop:p213-zero-current-smoothness-not-open-hereditary}]
Let $\pi:X\to P$ be the first projection.  It is a globally trivial
proper smooth curve bundle with fibre $D^{\mathrm{an}}$.  Hence
\Cref{prop:p062-proper-smooth-curve-bundle-descent} writes every continuous
zero-current function on $X$ uniquely as
\[
                              H=\pi^*b,
 \qquad b\in C^0(P,\mathbf R),
 \qquad \mathrm d'\mathrm d''[b]=0.
\]
By \Cref{cor:ch19-continuous-zero-current-harmonic}, the function $b$ is
classically pluriharmonic.  It attains a maximum because $P$ is compact.
The local maximum assertion in \Cref{prop:classical-psh-calculus} makes the
maximum locus open; it is also nonempty and closed.  Since $P$ is connected,
$b$, and therefore $H$, is constant.  This proves the first assertion.

Write $\mathcal H_D$ for the sheaf of Thuillier harmonic germs and
$\mathcal F_D$ for the real vector subsheaf generated by germs $\log|f|$
with $f$ an invertible analytic function.  By
\Cref{rem:ch08-harmonic-vs-log},
\[
       (\mathcal H_D/\mathcal F_D)_x
       \simeq\operatorname{Pic}^0(D_x)\otimes_{\mathbf Z}\mathbf R.
\]
Choose a nonzero class using
\eqref{eq:p213-nonhereditary-residual-class}, and represent it by a
Thuillier harmonic function $g$ on an ordinary open neighbourhood $U$ of
$x$.  Both $g$ and $-g$ are Thuillier-subharmonic.  The equivalence in
\Cref{thm:wanner-curve-comparison} therefore gives
$\mathrm d'\mathrm d''[g]=0$.  The residual Picard criterion
\Cref{prop:p213-residual-picard-smooth-curve} shows that $g$ is not
ordinary CLD-smooth near $x$.  Moreover,
\Cref{cor:ch19-continuous-zero-current-harmonic} makes $g$ harmonic in the
sense of Gubler--Jell--Rabinoff.

Let $q:\Omega=P\times_KU\to U$ be the second projection.  Harmonic
functoriality \Cref{prop:ch18-harmonic-functorial} makes
$G=q^*g$ harmonic.  The space $\Omega$ is good, topologically separated,
strict, and boundaryless.  Therefore
\cite[Corollary~7.15]{GJR26} gives
$\mathrm d'\mathrm d''[G]=0$.  Fix $p\in P(K)$ and let
\[
                         s_p:U\longrightarrow\Omega,
 \qquad z\longmapsto(p,z)
\]
be the analytic section.  If $G$ were ordinary CLD-smooth near $(p,x)$,
then functorial pullback of ordinary CLD forms in
\Cref{prop:cld-form-pullback} would make
$s_p^*G=g$ ordinary CLD-smooth near $x$, a contradiction.  This proves the
second assertion and the claimed failure of open inheritance.  Finally,
every type~$2$ residue curve of $P$ is rational, so $P$ satisfies the
residual-Picard condition.  The projection $X\to P$ is a locally trivial
proper smooth curve bundle, namely the fourth exact stage in
\Cref{prop:p062-mixed-exact-bundle-tower-transfer}; it is therefore one of
the exact stages allowed in the cited tower theorem.  This proves the last
scope assertion.
\end{proof}

\begin{proof}[of
\Cref{cor:p213-good-reduction-elliptic-open-nonheredity}]
The field $K$ is algebraically closed and complete, with residue field
$\mathbf C$.  The constant smooth proper model of $D$ has special fibre
$E$.  Its generic point defines a type~$2$ analytic point $x$, and the
associated smooth projective residue curve is $D_x\simeq E$.  The group
$E(\mathbf C)\simeq\operatorname{Pic}^0(E)$ contains a point of infinite
order: its torsion subgroup is countable, whereas $E(\mathbf C)$ is
uncountable.  Such a point remains nonzero after tensoring with
$\mathbf R$, so \eqref{eq:p213-nonhereditary-residual-class} holds.  Apply
\Cref{prop:p213-zero-current-smoothness-not-open-hereditary}.
\end{proof}
